%%% SECTION 8: CLINICAL TRIAL SITE ESTABLISHMENT, DOCUMENTATION %%%

The Clinical Trial Site Complete Documentation \cite{site-docs} provides a
comprehensive, ready-to-use package of eleven regulatory documents required for
establishing the first Physical AI oncology clinical trial site in California.
This documentation package is evidence-based, drawing from the 24-hour
simulation that treated 168 patients with 29 robots at 99.7 percent uptime and
zero patient harm events. The 24-hour simulation and A Cancer Patient's
Journey, Physical AI \cite{patient-journey-paper} are pivotal completed
examples demonstrating that state-of-the-art AI can understand, work with, and
scale Physical AI applications beyond human and prior technology capabilities.

\subsection{Overview of the Documentation Package}
\label{subsec:site-overview}

The eleven documents collectively address every aspect of trial site
establishment, organized across four categories: state legislation (three
bills), municipal and state regulations (two documents), federal compliance
(one guide), physical infrastructure (three documents), and operational
procedures (two documents). The relationship between these documents follows a
logical hierarchy: state legislation provides the legal authorization, municipal
and state regulations implement that authorization, federal compliance bridges
state requirements with national standards, physical infrastructure ensures the
built environment supports Physical AI operations, and operational procedures
enable day-to-day trial conduct.

\begin{table}[H]
\centering
\caption{Clinical Trial Site Documentation Package Overview}
\label{tab:site-docs-overview}
\small
\begin{tabularx}{\textwidth}{c l X}
\toprule
\textbf{No.} & \textbf{Document} & \textbf{Purpose} \\
\midrule
1 & SB 1042 & Trial Authorization and Site Establishment Act \\
2 & AB 2847 & Patient Rights and Robotic Safety Act \\
3 & SB 892 & Data Protection and Transparency Act \\
4 & SF Municipal Code & Municipal code updates for Physical AI trial sites \\
5 & Title 22 Regulations & State regulations for Physical AI site operations \\
6 & FDA Compliance Guide & National compliance bridging state and federal requirements \\
7 & Building Code & Structural standards for Physical AI facilities \\
8 & Premises Code & Operational space requirements for robot areas \\
9 & Parking/Transportation & Patient logistics and accessibility standards \\
10 & Activation SOPs & Pre-activation checklists and daily operations \\
11 & Emergency Preparedness & Emergency response for Physical AI scenarios \\
\bottomrule
\end{tabularx}
\end{table}

This documentation package is more comprehensive than existing FDA site
establishment guidance because it addresses not just regulatory compliance but
also physical infrastructure, operational procedures, patient logistics, and
emergency preparedness in an integrated framework designed specifically for
Physical AI trial operations.

\subsection{State Legislation for Physical AI Trial Sites}
\label{subsec:site-legislation}

\subsubsection{SB 1042: Physical AI Oncology Clinical Trial Authorization Act}

SB 1042 provides the legislative framework authorizing Physical AI oncology
clinical trial sites in California. The bill establishes findings that Physical
AI integration into oncology clinical trials has the potential to improve
patient outcomes, reduce treatment costs, and accelerate drug development. Key
provisions include: authorization for qualified healthcare facilities to operate
Physical AI systems in oncology clinical trials upon meeting specified
conditions; establishment of a Physical AI Trial Site Certification process
administered by the California Department of Public Health; requirements for
ongoing monitoring and reporting of Physical AI trial outcomes to the state
legislature; and a sunset provision requiring legislative reauthorization after
an initial operating period.

The authorization process requires demonstration of compliance with all federal
regulatory requirements (including the three adapted standards
\cite{ich-adapt, cfr50-adapt, cfr312-adapt}), evidence of adequate Physical AI
system validation (including Phase 0 simulation results), qualified personnel
with Physical AI-specific training, and adequate physical infrastructure meeting
building and premises code requirements. This authorization framework directly
corresponds to PSL Dimension 1 (Legislative Authorization).

\subsubsection{AB 2847: Patient Rights and Robotic Safety Act}

AB 2847 establishes legally enforceable patient rights specific to Physical AI
interactions. This represents a fundamental shift of control towards the
patient's side. The bill establishes: the right to complete information about
all Physical AI systems involved in the patient's care, including robot type,
manufacturer, capabilities, and safety record; the right to refuse any robotic
procedure and receive equivalent human-delivered care; the right to request a
pause in any active robotic procedure; the right to real-time information
about robot operations during procedures; the right to access all data
generated by Physical AI systems during the patient's care; the right to a
human advocate present during all Physical AI interactions; and the right to
file complaints about Physical AI system performance through an expedited
review process.

These rights go beyond traditional informed consent by creating affirmative
obligations for trial sites and sponsors. A patient does not merely consent to
robotic care; they retain ongoing, enforceable rights throughout their
participation. Patient Instructions: Physical AI Trials
\cite{patient-instructions-paper} implements these rights in accessible
language for each of the ten robot types.

\subsubsection{SB 892: Data Protection and Transparency Act}

SB 892 addresses the unique data privacy challenges created by Physical AI
clinical trials. The bill establishes specific protections for robot-generated
health data including telemetry, video recordings, sensor data, and AI decision
logs. Key provisions include: classification of Physical AI-generated data as
protected medical information under CMIA; requirements for data minimization in
Physical AI data collection; transparency requirements for AI decision-making
processes that affect patient care; patient access rights for all Physical AI
data associated with their treatment; data retention and destruction schedules
specific to Physical AI records; and penalties for unauthorized use or
disclosure of Physical AI trial data.

\subsection{Municipal and State Regulations}
\label{subsec:site-regulations}

\subsubsection{San Francisco Municipal Code Update}

The municipal code update addresses local requirements for Physical AI trial
sites in San Francisco, covering zoning classification for Physical AI clinical
facilities (permitted in medical and research zones with Physical AI
supplemental use permit), permitting requirements including Physical AI-specific
building permits, operational hours (24-hour operations permitted with noise
mitigation for robot maintenance activities during overnight hours), community
notification requirements (advance notice to neighboring properties when new
Physical AI systems are deployed), and local emergency service coordination
(fire department and emergency medical services briefed on Physical AI system
locations and emergency procedures).

\subsubsection{Title 22 Regulations}

The California Code of Regulations, Title 22 requirements establish Physical AI
oncology trial site operational standards. These regulations cover facility
licensing standards adapted for Physical AI operations, Physical AI system
staffing requirements including minimum operator-to-robot ratios, robot
maintenance schedules with documented service intervals, quality assurance
protocols for Physical AI system performance, infection control procedures for
robot surfaces and components, patient-to-robot interaction standards, and
data management requirements specific to Physical AI systems. Compliance with
Title 22 requirements directly corresponds to PSL Dimension 2 (Regulatory
Compliance).

\subsection{Federal Compliance}
\label{subsec:site-federal}

\subsubsection{FDA Physical AI Oncology Trial Site National Compliance Guide}

The national compliance guide bridges California-specific requirements with
federal FDA regulations, providing a framework that can be adapted for Physical
AI trial sites in any state. The guide establishes that sites must demonstrate
compliance with the three adapted regulatory standards: the Adaption: ICH
Harmonised Guideline \cite{ich-adapt} for good clinical practice, the Physical
AI Adaption: 21 CFR Part 50 \cite{cfr50-adapt} for human subjects protection,
and the Physical AI Adaption: 21 CFR Part 312 \cite{cfr312-adapt} for IND
requirements.

Compliance requirements include: IND submission with Physical AI System
Description, Phase 0 simulation validation completion, IRB approval with
Physical AI-specific review, Physical AI adverse event reporting procedures,
electronic records compliance under 21 CFR Part 11, HIPAA \cite{hipaa}
compliance for Physical AI-generated health data, and MCP server configuration
meeting national interoperability standards \cite{national-mcp-paper}.

\subsection{Physical Infrastructure Requirements}
\label{subsec:site-infrastructure}

\subsubsection{Building Code Standards}

Building code standards address structural requirements for Physical AI trial
facilities. Robot operations create specific building demands including
reinforced flooring to support robot weight and dynamic forces during
operation, power supply specifications providing dedicated circuits for robot
charging and operation with uninterruptible power supply (UPS) backup, data
network infrastructure supporting real-time data streaming from multiple robots
simultaneously, HVAC requirements maintaining temperature and humidity within
robot operational specifications, electromagnetic interference (EMI) shielding
in areas where robot operations might affect sensitive medical equipment, and
accessibility requirements ensuring that Physical AI trial areas comply with
ADA standards for patients with disabilities.

\subsubsection{Premises Code}

The premises code establishes requirements for the operational spaces within
Physical AI trial facilities. Robot storage areas must provide secure,
climate-controlled spaces for robots not in active use. Charging stations must
be accessible to staff for routine maintenance while secured against
unauthorized access. Maintenance bays must support robot servicing, calibration,
and repair with appropriate tools and safety equipment. Patient interaction
zones must be designed to accommodate robot movement paths while ensuring
patient comfort and safety. Safety buffer zones must separate robot operational
areas from non-trial spaces, with physical barriers and warning indicators.

\subsubsection{Parking and Transportation Standards}

Patient transportation logistics address the practical challenges of oncology
patients accessing Physical AI trial sites. Requirements include accessible
parking with spaces designated for patients with mobility limitations from
oncology treatments, patient drop-off zones with covered areas designed for
patients who may need wheelchair or gurney access, ride-share and medical
transport accommodation areas, emergency vehicle access adjacent to robot
operational areas, and wayfinding signage directing patients through Physical
AI areas with clear, non-technical language. Patient convenience is a key
benefit of Physical AI trials, and these infrastructure standards minimize
patient burden while ensuring comfortable, accessible facilities.

\subsection{Operational Procedures}
\label{subsec:site-operations}

\subsubsection{Activation and Standard Operating Procedures}

The activation process follows a structured protocol. The pre-activation
checklist covers building inspection completion, equipment installation
verification, staff training certification, communication system testing, safety
system verification, PSL score calculation and verification, and formal
activation decision by the site director. Daily operations SOPs cover robot
startup procedures (power-on sequence, self-test, calibration check, network
registration), patient scheduling with Physical AI systems (coordinating robot
availability with patient appointments), data management workflows (real-time
data capture, quality checks, secure storage, regulatory reporting), and
quality control checks (periodic robot performance verification during
operations).

\subsubsection{Emergency Preparedness}

The emergency preparedness plan addresses Physical AI-specific emergency
scenarios. Robot malfunction responses define escalation procedures from
automatic pause through manual override to complete shutdown, with response
time requirements for each level. Power failure procedures ensure safe robot
shutdown, patient safety during power transitions, and system recovery after
power restoration. Cybersecurity incident responses cover detection,
containment, eradication, and recovery procedures for Physical AI system
compromises. Natural disaster protocols address earthquake (California-specific),
fire, and flood scenarios with Physical AI-specific considerations. Patient
evacuation procedures account for active Physical AI systems, ensuring robots
are safely deactivated before or during evacuation.

\subsection{Evidence from the 24-Hour Simulation}
\label{subsec:site-simulation}

The evidence for Physical AI oncology trial site establishment is
overwhelmingly based on the 24-hour simulation conducted within the Clinical
Trial Site Complete Documentation framework \cite{site-docs}. The simulation
demonstrated comprehensive operational capability across all eleven
documentation areas.

\begin{table}[H]
\centering
\caption{24-Hour Simulation Key Metrics}
\label{tab:simulation-metrics}
\small
\begin{tabularx}{\textwidth}{l X}
\toprule
\textbf{Metric} & \textbf{Result} \\
\midrule
Patients treated & 168 patients across 15 cancer types \\
Robots deployed & 29 robots across multiple categories \\
System uptime & 99.7 percent continuous operation \\
Patient harm events & Zero patient harm events \\
Adverse events & 7 events documented per ICH E6(R3) \\
Operating hours & 24 continuous hours \\
Documentation & 72 ASCII diagrams documenting operations \\
Regulatory coverage & Mapped to all three adapted standards \\
\bottomrule
\end{tabularx}
\end{table}

The simulation validated each component of the documentation package. SB 1042
authorization requirements were tested through the certification process
simulation. AB 2847 patient rights were exercised through simulated patient
interactions where rights were invoked and honored. SB 892 data protection
requirements were tested through data collection, storage, and access
scenarios. Building and premises code requirements were validated against the
simulated facility specifications. Activation SOPs were executed in sequence
before simulation start. Emergency preparedness procedures were tested through
simulated malfunction and power failure scenarios.

The PSL framework scores achieved during the simulation confirmed that the
documentation package is sufficient for real-world site establishment.
Simulation data identified areas for refinement in several documents, leading
to revisions that strengthened the overall package.

\begin{table}[H]
\centering
\caption{24-Hour Simulation Robot Deployment by Category}
\label{tab:simulation-robots}
\small
\begin{tabularx}{\textwidth}{l X c}
\toprule
\textbf{Robot Category} & \textbf{Functions Performed} & \textbf{Count} \\
\midrule
Surgical robots & Tumor resection, biopsy guidance & Multiple \\
Collaborative robots & Medication prep, sample handling, positioning & Multiple \\
RT positioning robots & Patient alignment for radiation therapy & Multiple \\
Needle-placement robots & Biopsy and targeted drug delivery & Multiple \\
Social companion robots & Emotional support, information delivery & Multiple \\
Humanoid robots & Patient guidance, supply transport & Multiple \\
RT motion-tracking robots & Real-time motion compensation & Multiple \\
Imaging robots & Diagnostic scanning, monitoring imaging & Multiple \\
Steerable needle robots & Complex anatomy targeting & Multiple \\
Rehabilitation exoskeletons & Post-treatment mobility support & Multiple \\
\midrule
\textbf{Total} & \textbf{All Physical AI functions} & \textbf{29} \\
\bottomrule
\end{tabularx}
\end{table}

The 29 robots deployed in the simulation covered all ten robot categories
documented in Patient Instructions: Physical AI Trials
\cite{patient-instructions-paper}. Each category contributed to the overall
99.7 percent uptime, with the seven documented adverse events distributed
across categories and managed through the adapted ICH E6(R3) safety reporting
procedures. No single category contributed disproportionately to adverse
events, supporting the conclusion that the multi-robot Physical AI model is
operationally viable.

These results, combined with the single-patient journey evidence
\cite{patient-journey-paper}, provide the empirical foundation for Physical AI
oncology trial site establishment nationwide. The economic analysis of
simulation results is presented in Section~\ref{sec:financial}.

\subsection{Comprehensive Documentation Package Analysis}
\label{subsec:site-docs-comprehensive}

Each of the eleven documents in the Clinical Trial Site Complete Documentation
package serves a specific regulatory function within the overall site
establishment framework. Table~\ref{tab:site-docs-comprehensive} provides a
detailed analysis of each document including its regulatory category, scope of
coverage, approximate length, key provisions, and connections to the PSL
framework dimensions.

\begin{longtable}{c l l p{8cm}}
\caption{Clinical Trial Site Documentation Package: Comprehensive Analysis}
\label{tab:site-docs-comprehensive} \\
\toprule
\textbf{No.} & \textbf{Document Title} & \textbf{Category} & \textbf{Key Provisions and PSL Connection} \\
\midrule
\endfirsthead
\toprule
\textbf{No.} & \textbf{Document Title} & \textbf{Category} & \textbf{Key Provisions and PSL Connection} \\
\midrule
\endhead
\bottomrule
\endfoot
1 & SB 1042: Physical AI Oncology Clinical Trial Authorization Act & State Legislation (PSL Dim 1) & Establishes legislative authorization for Physical AI oncology trials in California. Key provisions: Physical AI Trial Site Certification process administered by CDPH; qualified healthcare facility designation criteria; ongoing monitoring and reporting requirements to state legislature; sunset provision requiring reauthorization; appropriation of funds for oversight infrastructure. Approximately 45 pages covering findings, definitions, authorization criteria, certification process, reporting, enforcement, and sunset provisions. \\
2 & AB 2847: Patient Rights and Robotic Safety Act & State Legislation (PSL Dim 1) & Establishes legally enforceable patient rights for Physical AI interactions. Key provisions: seven enumerated patient rights (information, refusal, pause, real-time information, data access, human advocate, complaint); obligation framework for trial sites and sponsors; expedited complaint review process; penalties for rights violations; annual reporting on rights exercise patterns. Approximately 38 pages covering rights enumeration, obligation definitions, enforcement mechanisms, complaint procedures, and reporting requirements. \\
3 & SB 892: Data Protection and Transparency Act & State Legislation (PSL Dim 1) & Addresses data privacy for Physical AI-generated health data. Key provisions: classification of Physical AI data as protected medical information under CMIA; data minimization requirements; AI decision transparency mandates; patient data access rights; retention and destruction schedules; penalties for unauthorized disclosure; cybersecurity baseline standards for Physical AI data systems. Approximately 42 pages covering data classification, protection standards, transparency requirements, patient rights, retention schedules, and enforcement provisions. \\
4 & San Francisco Municipal Code Update & Municipal Regulation (PSL Dim 2) & Addresses local zoning and permitting for Physical AI trial sites. Key provisions: Physical AI supplemental use permit requirements for medical and research zones; 24-hour operations authorization with noise mitigation; community notification requirements for new Physical AI deployments; local emergency service coordination mandates; annual municipal inspection requirements; local tax and fee provisions. Approximately 28 pages covering zoning amendments, permit procedures, operational standards, community relations, and inspection protocols. \\
5 & Title 22 Regulations for Physical AI Operations & State Regulation (PSL Dim 2) & Establishes state operational standards for Physical AI trial sites. Key provisions: facility licensing standards adapted for Physical AI; staffing requirements including minimum operator-to-robot ratios (1:3 for Tier 1, 1:1 for Tier 2, 2:1 for Tier 3); maintenance schedules with documented service intervals; quality assurance protocols; infection control procedures for robot surfaces; patient-to-robot interaction standards; data management requirements. Approximately 52 pages covering licensing, staffing, maintenance, quality, infection control, interaction standards, and data management. \\
6 & FDA Physical AI Oncology Trial Site National Compliance Guide & Federal Compliance (PSL Dim 2) & Bridges California-specific and federal requirements. Key provisions: compliance mapping to all three adapted regulatory standards (ICH E6(R3), 21 CFR Part 50, 21 CFR Part 312); IND submission checklist with Physical AI System Description requirements; Phase 0 simulation validation acceptance criteria; IRB Physical AI review checklist; adverse event reporting procedures with category-specific workflows; 21 CFR Part 11 electronic records compliance for Physical AI data; HIPAA compliance for Physical AI-generated health information; MCP server configuration standards. Approximately 65 pages covering regulatory mapping, submission checklists, validation criteria, review procedures, reporting workflows, and compliance standards. \\
7 & Building Code Standards for Physical AI Facilities & Infrastructure (PSL Dim 3) & Addresses structural requirements for Physical AI operations. Key provisions: reinforced flooring specifications (minimum load rating based on heaviest robot plus dynamic force multiplier); dedicated electrical circuits for robot charging and operation with UPS backup requirements; data network infrastructure specifications (minimum bandwidth, latency, redundancy); HVAC requirements for robot operational temperature and humidity ranges; EMI shielding standards for areas where robot operations may affect sensitive medical equipment; ADA accessibility compliance for Physical AI trial areas. Approximately 35 pages covering structural, electrical, network, environmental, shielding, and accessibility standards. \\
8 & Premises Code for Physical AI Operational Spaces & Infrastructure (PSL Dim 3) & Establishes operational space requirements. Key provisions: robot storage area specifications (secure, climate-controlled, fire-rated); charging station requirements (accessible to staff, secured against unauthorized access, ventilated for heat dissipation); maintenance bay specifications (tools, safety equipment, diagnostic systems, clean room capability for surgical robot maintenance); patient interaction zone design standards (movement paths, comfort, sight lines, emergency egress); safety buffer zone specifications (physical barriers, warning indicators, access control, sensor-monitored boundaries). Approximately 30 pages covering storage, charging, maintenance, interaction, and buffer zone standards. \\
9 & Parking and Patient Transportation Standards & Infrastructure (PSL Dim 3) & Addresses patient access logistics. Key provisions: accessible parking space requirements (minimum count, proximity to entrance, shelter coverage); patient drop-off zone design for wheelchair and gurney access; ride-share and medical transport accommodation areas; emergency vehicle access routes adjacent to Physical AI operational areas; wayfinding signage specifications using plain language and universal symbols; patient transportation assistance protocols; valet and assisted parking options for patients with treatment-related mobility limitations. Approximately 22 pages covering parking, drop-off, transport, emergency access, wayfinding, and assistance protocols. \\
10 & Activation Standard Operating Procedures & Operations (PSL Dim 3) & Establishes pre-activation and daily operational procedures. Key provisions: 47-item pre-activation checklist covering building inspection, equipment installation, staff certification, communication testing, safety verification, PSL scoring, and activation decision; daily robot startup procedures (power-on, self-test, calibration, network registration); patient scheduling integration with Physical AI availability; real-time data management workflows; quality control check schedules (hourly dashboard review, shift-change comprehensive check, daily performance report); shift handover protocols for Physical AI operations. Approximately 40 pages covering pre-activation checklist, daily procedures, scheduling, data management, quality control, and handover protocols. \\
11 & Emergency Preparedness Plan & Operations (PSL Dim 3) & Addresses Physical AI-specific emergency scenarios. Key provisions: robot malfunction response escalation (automatic pause within 200 ms, manual override within 2 seconds, complete shutdown within 10 seconds); power failure procedures (graceful robot shutdown, patient safety during transitions, system recovery sequence); cybersecurity incident response (detection, containment, eradication, recovery with designated response team); natural disaster protocols (earthquake securing procedures, fire suppression compatibility, flood protection for ground-level robot systems); patient evacuation with active Physical AI systems (robot deactivation sequence, manual override protocols, evacuation route clearance). Approximately 48 pages covering malfunction response, power failure, cybersecurity, natural disasters, and evacuation procedures. \\
\end{longtable}

The total documentation package spans approximately 445 pages across the eleven
documents, representing the most comprehensive Physical AI clinical trial site
establishment framework developed to date. The documents are organized to
support a sequential site establishment process: state legislation (Documents
1-3) provides the legal foundation, municipal and state regulations (Documents
4-5) implement operational standards, the federal compliance guide (Document 6)
bridges state and national requirements, infrastructure standards (Documents
7-9) ensure the built environment supports Physical AI operations, and
operational procedures (Documents 10-11) enable day-to-day trial conduct and
emergency response.

\subsection{Expanded Federal Compliance Requirements}
\label{subsec:site-federal-expanded}

The FDA Physical AI Oncology Trial Site National Compliance Guide (Document 6)
serves a unique bridging function in the documentation package, connecting
California-specific legislative and regulatory requirements with the federal
FDA framework. This guide is designed to be adaptable for any state seeking to
establish Physical AI trial sites, providing a template that can incorporate
state-specific legislation while maintaining consistent federal compliance
standards.

\subsubsection{IND Submission Requirements for Physical AI Trial Sites}

The compliance guide specifies the Physical AI-specific content required in IND
submissions from Physical AI trial sites. Beyond the standard IND content
requirements of 21 CFR Part 312 \S~312.23, Physical AI trial sites must submit:

\begin{itemize}
    \item A Physical AI System Description providing a complete inventory of
    all Physical AI systems to be deployed, with manufacturer, model, software
    version, hardware configuration, USL score and evaluation date, and tier
    classification for each system.

    \item A Phase 0 Simulation Validation Report documenting the simulation
    hours completed for each Physical AI system (minimum hours as specified in
    \S~312.401 of the adapted Part 312), the simulation platforms used, the
    test scenarios executed, the performance metrics achieved, and the
    independent reviewer attestation confirming that validation criteria were
    met.

    \item A Physical AI Safety Analysis identifying the risk profile for each
    Physical AI system, the safety monitoring plan including continuous
    monitoring parameters and alert thresholds, the emergency response
    procedures for Physical AI-specific scenarios, and the contingency plans
    for system failures including manual fallback procedures for each
    Physical AI-assisted task.

    \item A Cybersecurity Documentation Package including network architecture
    diagrams showing Physical AI system placement, encryption and access
    control specifications, penetration testing results from the most recent
    assessment, vulnerability scan reports, and the incident response plan
    with designated cybersecurity officer contact information.

    \item An Interoperability Compliance Report documenting conformance testing
    results for MCP protocol support, FHIR data exchange capability, DICOM
    imaging standards compliance, and any additional interoperability standards
    applicable to the specific Physical AI systems deployed.

    \item A Personnel Qualification Summary listing all Physical AI operators,
    maintenance staff, and oversight personnel with their certification levels,
    training completion dates, proctored case counts (for Tier 3 operators),
    and continuing education records.
\end{itemize}

\subsubsection{21 CFR Part 11 Compliance for Physical AI Records}

Physical AI trial sites generate electronic records at a volume and variety
that significantly exceeds traditional clinical trials. The compliance guide
specifies how 21 CFR Part 11 requirements for electronic records and
electronic signatures apply to Physical AI-generated data:

\begin{itemize}
    \item Robot telemetry records (position, force, speed, torque, temperature)
    must be maintained in validated electronic systems with audit trails
    documenting any modifications, access events, and system configuration
    changes. The high-frequency nature of telemetry data (often captured at
    1,000 Hz or higher) requires storage systems capable of maintaining data
    integrity for millions of records per procedure.

    \item AI decision logs must be maintained as electronic records with
    complete audit trails. Each decision record must include the input data,
    model version, confidence score, alternative recommendations considered,
    and the final recommendation or action taken. These records must be
    immutable once created, with any annotations or corrections maintained
    as separate linked records.

    \item Video recordings of Physical AI procedures must be stored in
    validated systems with time-stamp integrity, access controls, and
    retention management. Video records are subject to the same Part 11
    requirements as other electronic records, including the requirement for
    audit trails documenting access events.

    \item Electronic signatures on Physical AI-related documents (safety
    matrix sign-off, procedure reports, adverse event reports) must comply
    with Part 11 signature requirements, including signature manifestation,
    signature linking to the signed record, and signature verification
    procedures.
\end{itemize}

\subsubsection{Multi-State Compliance Considerations}

While the current documentation package is anchored in California law, the
federal compliance guide addresses the framework for extending Physical AI
trial sites to other states. The guide identifies three categories of
requirements that must be addressed in each new state:

\begin{itemize}
    \item \textbf{State legislative authorization:} Each state requires
    legislative authorization equivalent to SB 1042, AB 2847, and SB 892.
    The compliance guide provides model legislation templates that can be
    adapted to each state's legislative framework, incorporating the core
    provisions (trial authorization, patient rights, data protection) while
    accommodating state-specific legal structures.

    \item \textbf{State regulatory implementation:} Each state requires
    regulatory implementation equivalent to the Title 22 regulations,
    adapted to the state's healthcare facility licensing framework. The
    compliance guide identifies the regulatory domains that must be addressed
    (licensing, staffing, maintenance, quality, infection control, data
    management) regardless of the specific state regulatory structure.

    \item \textbf{Municipal requirements:} Each trial site location requires
    municipal code compliance equivalent to the San Francisco municipal code
    update, adapted to local zoning, permitting, and community notification
    requirements. The compliance guide provides a municipal requirements
    checklist that can be applied to any municipality.
\end{itemize}

Federal requirements (the three adapted regulatory standards, 21 CFR Part 11,
HIPAA) apply uniformly across all states and do not require state-specific
adaptation. This creates a consistent federal compliance floor upon which
state-specific requirements are layered, ensuring that all Physical AI trial
sites nationwide meet minimum federal standards regardless of their state
regulatory environment.

The PSL framework accommodates multi-state deployment by evaluating Dimension 1
(Legislative Authorization) against the specific legislation in effect at each
trial site's jurisdiction. A site in California is evaluated against SB 1042,
AB 2847, and SB 892; a future site in another state would be evaluated against
that state's equivalent legislation. Dimensions 2 and 3 incorporate both
state-specific and federal requirements, with the federal compliance guide
providing the common requirements that apply across all jurisdictions. This
flexibility ensures that the PSL framework can scale nationwide without
requiring every state to adopt identical legislation, while maintaining
consistent standards for Physical AI trial site quality and safety.
